\scanpage{201}
% archive sheet number: 200
1) \uline{Cohomologie des diviseurs à croisements normaux.}

\[
\begin{aligned}
\underset{\text{propre}}{Y}=\bigcup_\alpha Y_\alpha\qquad &Y_\alpha\ \text{projectif et lisse,}\\
&Y_{\alpha_0}\cap\cdots\cap Y_{\alpha_p}\ \text{est lisse.}
\end{aligned}
\]
\[
H^{\unsure{*|\bullet}}(Y,\mathbb{Q}_{\unsure{\ell|l}})\Longleftarrow E_2^{pq}=H^p(\sigma\mapsto H^q(Y_\sigma,\mathbb{Q}_\ell)),\qquad
Y_{\alpha_0\cdots\alpha_p}=Y_{\alpha_0}\cap\cdots\cap Y_{\alpha_p}
\]
\deleted{\unsure{Proposition}} \deleted{\unsure{Théorème}}

$E_2^{pq}$ est de \uline{poids} $q$ donc $d_r^{pq}$ est nulle pour tout $r$.
% slanted
\margin{\uline{NB} de plus (\unsure{via} conjectures standard) les dimensions des $E_2^{pq}$ sont indépendantes de $\ell$.}
$H^i(Y,\mathbb{Q}_\ell)$ \unsure{se dévisse} en les $E_2^{p,q}$ $(p+q=i)$, de poids $0,1,\ldots,i$.

2) Cohomologie des \deleted{variétés} \inserted{schémas} quasi-projectifs lisses non propres.

$X=X'-Y$, $X'$ projectif et lisse, $Y$ diviseur à croisements normaux, tels que les comp. irréduc. $Y_i$ soient lisses.
\[
\begin{aligned}
H^{i-1}(X',\mathbb{Q}_\ell)\to\underset{\text{poids }\leq i-1}{H^{i-1}(Y,\mathbb{Q}_\ell)}\to H^i_{\deleted{\unsure{c}}!}(X,\mathbb{Q}_\ell)\to\underset{\text{poids }i}{H^i(X',\mathbb{Q}_\ell)}\to\underset{\text{poids }\leq i}{H^i(Y,\mathbb{Q}_\ell)}\to H^{i+1}_!(X,\mathbb{Q}_\ell)\to\cdots
\end{aligned}
\tag*{\textcircled{1}}
\]

On trouve que
\[
\text{a) }\deleted{\text{l'image}}\ \operatorname{Im}(H^{i-1}(Y,\mathbb{Q}_\ell)\inserted{\to H^i_!(X,\mathbb{Q}_\ell)})=\operatorname{Ker}(H^i_!(X,\mathbb{Q}_\ell)\unsure{-|\to}H^i(X',\mathbb{Q}_\ell))
\]
\uline{ne dépend pas} de la résolution choisie : c'est le morceau de poids $\leq i-1$ de $H^i_!(X,\mathbb{Q}_\ell)$.

\deleted{De plus, le \unsure{Ker et le Coker}}

\unsure{Le|Les} Ker et Coker de $H^i(X',\mathbb{Q}_\ell)\overset{\varphi^i}{\unsure{-|\to}}H^i(Y,\mathbb{Q}_\ell)$, à isom. can. près, est aussi indépendant de la résolution choisie. \deleted{\unsure{Ce dernier se dévisse en un}} \inserted{Considérons la filtration de} $H^i(Y,\mathbb{Q}_\ell)$ provenant de la suite spectrale étudiée dans 1), correspond. à un \unsure{déviss. gr} en $E_2^{pq}=H^p(\sigma\mapsto H^q(Y_\sigma,\mathbb{Q}_\ell))$ $(p+q=i)$. L'hom.
% formula crowded against the right and lower edges
\[
H^i(Y,\mathbb{Q}_\ell)\unsure{-|\to}E_2^{0,i}=H^0(\sigma\mapsto H^i(Y_\sigma,\mathbb{Q}_\ell))=\operatorname{Ker}(\deleted{\unsure{\bigoplus}}\unsure{\prod_\alpha}H^i(Y_\alpha,\mathbb{Q}_\ell)\inserted{\to\cdots})
\]
